\uuid{8mf5}
\titre{Changement de variables}
\auteur{hanani ; mzari}
\organisation{exo7}
\licence{CC BY-NC-SA 4.0}
\datecreate{2020-01-20}
\niveau{L1}
\module{Analyse}
\chapitre{Calcul d'intégrales}
\sousChapitre{Primitives diverses}
\difficulte{3}
\theme{Primitives, changement de variables}
\exo7qcmid{lille-2.primitives.039}

\contenu{
\begin{question}
\qclassification{127.11}
\qtags{motcle=changement de variables}
On note par $F$ une primitive sur $\Rr$ de $\displaystyle f(x)=\frac{1}{2+4x^2-4x}$. Parmi les assertions suivantes, cocher celles qui sont vraies :
\begin{answers}
\good{$\displaystyle f(x)=\frac{1}{1+(2x-1)^2}$.}
\bad{$\displaystyle F(x)=\int \frac{\mathrm{d}u}{1+u^2}$ avec $u=2x-1$.}
\bad{$\displaystyle F(x)=\arctan (2x-1)+k$, $k\in \Rr$.}
\good{$\displaystyle F(x)=\frac{1}{2}\arctan (2x-1)+k$,\; $k\in \Rr$.}
\end{answers}
\begin{explanations}
On a : $2+4x^2-4x=1+(2x-1)^2$ et avec $u=2x-1$, on a : $\mathrm{d}u=2\mathrm{d}x$ et
$$F(x)=\frac{1}{2}\int \frac{\mathrm{d}u}{1+u^2}=\frac{1}{2}\arctan u+k=\frac{1}{2}\arctan (2x-1)+k,\; k\in \Rr.$$
\end{explanations}
\end{question}
}
